\chapter{Conclusion}
\label{chap:conclusion}

\section{Summary of Work}

This project designed, implemented, and evaluated a reproducible benchmarking
framework for Monte Carlo option-pricing engines with automatic
differentiation. The framework supports explicit workload configurations,
engine abstraction, deterministic seeding, warm-up and repeated timing,
analytical validation where available, SQLite result storage, environment
metadata capture, cloud cost metadata, and generated report artefacts.

The completed implementation compares NumPy, JAX/XLA, C++/OpenMP, and
Rust/Rayon engines across European GBM, parametric local-volatility, and Asian
option workloads where supported. The local experimental campaign provides
detailed evidence on validation, runtime, throughput, AD overhead,
path-count scaling, time-step scaling, thread scaling, memory pressure,
surface-shape sensitivity, and timing stability. The final cloud experiment
extends this with cost-performance analysis, Pareto-frontier analysis, and
profiler-versus-grid evaluation.

\section{Key Findings}

The results show that MC+AD performance cannot be reduced to a single backend
ranking. The simple European GBM workload is useful for correctness checks, but
the more demanding local-volatility and Asian workloads change the observed
trade-offs. Rust/Rayon is the strongest no-AD engine in the main
local-volatility results and in the highlighted European/local-volatility cloud
rows, while JAX/XLA remains essential for AD and is competitive for some
path-dependent configurations.

The AD results show a clear runtime-memory trade-off. Reverse-mode AD is
substantially cheaper than forward mode for the scalar-output
local-volatility workload at moderate sizes, but it reaches practical memory
limits at larger \(M,N\) combinations. In the final cloud comparison at
\(M=100{,}000\), JAX AD overhead ranges from \(1.81\times\) for Asian reverse
mode to \(4.42\times\) for local-volatility forward mode.

The cloud-profiler result is the strongest final systems contribution. In the
\code{sha\_cloud\_profiler\_v4} experiment on \code{n2-standard-4}, the full
grid requires 48 full benchmark cells. SHA selects 24, saving 50.0\% of full
runs while preserving 100.0\% Pareto recovery and 0.0\% runtime and cost regret
relative to the full grid. The old profiler also preserves the best decisions,
but selects 30 full runs, so SHA achieves the same measured decision quality
with fewer full-scale evaluations.

\section{Answering Research Questions}

\paragraph{Hardware and software stacks.}
The local and cloud results show that runtime depends strongly on workload
structure, engine implementation, and runtime system. Compiled CPU engines are
very strong for no-AD workloads, while JAX/XLA is necessary for the reported AD
measurements and can be competitive when the workload has enough structure to
benefit from compilation.

\paragraph{Programming languages and AD libraries.}
C++/OpenMP and Rust/Rayon provide high no-AD throughput and explicit CPU
parallelism. JAX provides the main AD capability and exposes the practical
difference between forward and reverse mode. The trade-off is therefore not
only language performance, but also derivative support, memory behaviour, and
ease of integration into a reproducible benchmark runner.

\paragraph{Cloud cost-performance.}
The final priced comparison shows that runtime and cost should be considered
together. For the measured no-AD rows at \(M=100{,}000\),
\code{t2d-standard-4} is faster and lower cost per run than
\code{n2-standard-4}. This is a measured result under the stored pricing and
environment assumptions, not a universal claim about all cloud hardware.

\paragraph{AI/ML techniques.}
JAX/XLA demonstrates that compiler and AD techniques from the AI/ML ecosystem
are relevant to MC workloads, especially when the computation is expressed in a
form such as \code{jax.lax.scan} that the compiler can optimise. The project
does not complete mixed precision, GPU auto-tuning, or broader hardware-aware
scheduling experiments, so the conclusion is deliberately limited to the
measured XLA and AD evidence.

\paragraph{Parallelism and bottlenecks.}
The thread-scaling experiments show effective scaling for C++ and Rust under
OpenMP/Rayon-style parallelism, while JAX's CPU runtime is less directly
controlled by requested thread counts. The stress tests show that memory, not
only arithmetic throughput, becomes a limiting factor for reverse-mode AD on
large local-volatility workloads.

\section{Future Work}

The most direct next step is to broaden the cloud study. More instance types,
regions, repeated cloud runs, and up-to-date pricing snapshots would test
whether the SHA profiler remains reliable on larger and noisier grids. The
profiler could also incorporate uncertainty estimates so that promotion
decisions account explicitly for timing variance and Monte Carlo noise.

GPU execution, distributed execution, mixed precision, and hardware counters
remain important extensions. The framework already has a CUDA scaffold and
metadata structure, but the final evaluation does not claim GPU results.
Extending the benchmark suite to GPUs and distributed workers would make the
hardware-comparison part of the original proposal much stronger.

Further AD systems and languages would also be valuable. The project evaluates
JAX AD, but future work could compare Enzyme, C++ AD tools, Rust AD libraries,
Mojo, or domain-specific quant libraries. Finally, the frontend and API could
be developed into a more complete dashboard for interactive profiler
configuration and report artefact generation.

\section{Final Remarks}

Overall, the project contributes both a working engineering artefact and an
empirical methodology. The artefact makes MC+AD experiments reproducible,
queryable, and extensible across engines. The methodology shows how detailed
local benchmarking can be combined with cloud cost metadata and successive
halving to reduce expensive full-grid benchmarking while preserving measured
selection quality. The final report therefore represents an expanded and
completed version of the original framework work: it retains the detailed
technical and experimental narrative, while adding the final cloud,
cost-performance, Pareto-frontier, AD-overhead, and SHA-profiler results.
